% Generated draft: requires mathematical and source review.
\documentclass[11pt]{article}
\usepackage[margin=25mm]{geometry}
\usepackage{fontspec}
\setmainfont{Latin Modern Roman}
\usepackage{amsmath,amssymb,mathtools,unicode-math}
\setmathfont{Latin Modern Math}
\usepackage{enumitem,needspace,xurl,hyperref}
\setlength{\parindent}{0pt}
\setlength{\parskip}{5pt}
\setlength{\emergencystretch}{3em}
\newcommand{\N}{\mathbb{N}}
\newcommand{\Z}{\mathbb{Z}}
\newcommand{\Q}{\mathbb{Q}}
\newcommand{\R}{\mathbb{R}}
\newcommand{\C}{\mathbb{C}}
\newcommand{\F}{\mathbb{F}}
\newcommand{\A}{\mathbb{A}}
\newcommand{\PP}{\mathbb{P}}
\newcommand{\E}{\mathbb{E}}
\newcommand{\eps}{\varepsilon}
\newcommand{\dd}{\,\mathrm{d}}
\begin{document}
\section*{UnsolvedMath ID: 800218}
\textbf{Source problem code:} catalogue entry
\subsection{Definitions and mathematical statement}
A biholomorphism between complex manifolds is a bijective holomorphic map with holomorphic inverse. Write \(X\cong_{\mathrm{bihol}}Y\) when such a map exists.

The supplied source calls a complex manifold \(X\) a \emph{long \(\C^2\)} if it is a union of proper subsets biholomorphic to \(\C^2\). Written literally, this condition is
\[
X=\bigcup_{i\in I}U_i,\qquad U_i\subsetneq X,\qquad U_i\cong_{\mathrm{bihol}}\C^2\quad(i\in I).
\]
Here the source does not specify the index set \(I\), any ordering or nesting of the subsets, or explicitly the complex-manifold structures on those subsets.

The question is: must every long \(\C^2\), in the intended sense of the source, satisfy
\[
X\cong_{\mathrm{bihol}}\C^2?
\]
The intended definition requires clarification: the supplied wording does not require a countable increasing exhaustion by open subsets, and such requirements must not be silently added. The displayed covering condition records the supplied wording; it does not settle these missing qualifications.
\subsection{Short English statement}
The source asks whether every complex manifold assembled as a union of proper copies of \(\C^2\) is itself biholomorphic to \(\C^2\). The permitted covering and the intended complex structures on its members are not fully specified.
\subsection{Sources}
\begin{itemize}
\item Local UnsolvedMath registry, permanent ID \(800218\), source code \texttt{catalogue entry}, title \emph{Several Complex Variables Problem List --- Problem 14}. The registry attributes the problem to Buzzard, Fornaess, and Krantz, with proposed year 1995, and labels it open; that label is supplied metadata, not an independently verified assessment.
\item Buzzard, Fornaess, and Krantz, \emph{The Several Complex Variables Problem List} (1995), Problem 14. Supplied source URL: \url{https://arxiv.org/abs/math/9509202}. The record reports access on 2026-10-16 and extraction by \texttt{arxiv\_latex}; these are supplied provenance statements, not a claim of independent source access.
\end{itemize}
\end{document}
